%-------------------------
% 中文简历 (Chinese Resume)
% 模板基于 Jake Gutierrez / sb2nov 简历模板改造
% 编译器: XeLaTeX (Overleaf: Menu -> Compiler -> XeLaTeX)
%------------------------

\documentclass[letterpaper,11pt]{article}

\usepackage[UTF8, fontset=fandol]{ctex}
% ctex 默认 1.3 倍行距，简历压缩为 1.08 以控制在两页内
\linespread{1.08}
\usepackage{latexsym}
\usepackage{makecell}
\usepackage[empty]{fullpage}
\usepackage{titlesec}
\usepackage[usenames,dvipsnames]{color}
\usepackage{verbatim}
\usepackage{enumitem}
\usepackage[hidelinks]{hyperref}
\usepackage{fancyhdr}
\usepackage{tabularx}

\pagestyle{fancy}
\fancyhf{}
\fancyfoot{}
\renewcommand{\headrulewidth}{0pt}
\renewcommand{\footrulewidth}{0pt}

% Adjust margins
\addtolength{\oddsidemargin}{-0.5in}
\addtolength{\evensidemargin}{-0.5in}
\addtolength{\textwidth}{1in}
\addtolength{\topmargin}{-.5in}
\addtolength{\textheight}{1.0in}

\urlstyle{same}
\raggedbottom
\raggedright
\setlength{\tabcolsep}{0in}

% Sections formatting
\titleformat{\section}{
  \vspace{-4pt}\raggedright\large\bfseries
}{}{0em}{}[\color{black}\titlerule \vspace{-5pt}]

%-------------------------
% Custom commands
\newcommand{\resumeItem}[1]{
  \item\small{%
    {#1\vspace{-2pt}}%
  }
}

\newcommand{\resumeSubheading}[4]{
  \vspace{-2pt}\item
    \begin{tabular*}{0.97\textwidth}[t]{l@{\extracolsep{\fill}}r}
      \textbf{#1} & #2 \\
      \textit{\small#3} & \textit{\small #4} \\
    \end{tabular*}\vspace{-7pt}
}

\newcommand{\resumeSubSubheading}[2]{
    \item
    \begin{tabular*}{0.97\textwidth}{l@{\extracolsep{\fill}}r}
      \textit{\small#1} & \textit{\small #2} \\
    \end{tabular*}\vspace{-7pt}
}

\newcommand{\resumeProjectHeading}[2]{
    \item
    \begin{tabular*}{0.97\textwidth}{l@{\extracolsep{\fill}}r}
      \small#1 & #2 \\
    \end{tabular*}\vspace{-7pt}
}

\newcommand{\resumeSubItem}[1]{\resumeItem{#1}\vspace{-4pt}}

\renewcommand\labelitemii{$\vcenter{\hbox{\tiny$\bullet$}}$}

\newcommand{\resumeSubHeadingListStart}{\begin{itemize}[leftmargin=0.15in, label={}]}
\newcommand{\resumeSubHeadingListEnd}{\end{itemize}}
\newcommand{\resumeItemListStart}{\begin{itemize}}
\newcommand{\resumeItemListEnd}{\end{itemize}\vspace{-5pt}}

%-------------------------------------------
\begin{document}
\raggedbottom

%----------HEADER----------
\begin{center}
    \textbf{\LARGE 张臻昊 (Zhenhao Zhang)} \\ \vspace{2pt}
    \small 纽约, NY $|$ (585) 470-3118 $|$ \href{mailto:zz3530@columbia.edu}{zz3530@columbia.edu} $|$
    \href{https://www.linkedin.com/in/zhenhao-zhang-uofr/}{LinkedIn} $|$
    \href{https://github.com/X0X0X00}{GitHub} \\ \vspace{2pt}
    \small AI 智能体 $\cdot$ 多模态学习 $\cdot$ 强化学习 $\cdot$ 计算机操作智能体 (Computer-Use Agents)
\end{center}

%-----------EDUCATION-----------
\section{教育背景}
\resumeSubHeadingListStart
  \resumeSubheading
    {哥伦比亚大学 (Columbia University)}{纽约, NY}
    {人工智能硕士 (M.S. in Artificial Intelligence)；方向：AI 与高级计算}{2026.08 -- 2027.12（预计）}
  \resumeSubheading
    {罗切斯特大学 (University of Rochester)}{罗切斯特, NY}
    {计算机科学理学学士；商学（信息系统）理学学士（双学位）}{2022.08 -- 2025.12}
  \resumeItemListStart
    \resumeItem{GPA：\textbf{3.85/4.00}；Cum Laude 荣誉毕业（计算机科学与商学信息系统）；计算机科学最高荣誉（Highest Distinction）；多学期 Dean's List（2023--2025）}
  \resumeItemListEnd
\resumeSubHeadingListEnd

%-----------PUBLICATIONS-----------
\section{论文发表}
\resumeSubHeadingListStart
  \resumeProjectHeading
    {\href{https://arxiv.org/abs/2609.31957}{\makecell[l]{CaptchaArena: A Large-Scale, Fine-Grained Dataset for Training Computer-Use Agents\\on Interactive CAPTCHAs}}}{\textbf{arXiv, 2026}}
  \resumeProjectHeading
    {\href{https://arxiv.org/abs/2602.03664}{Mitigating Conversational Inertia in Multi-Turn Agents}}{\textbf{ICML 2026}}
  \resumeProjectHeading
    {\href{https://arxiv.org/abs/2607.23263}{SeekJudge: A Practical Reward Framework for Reinforcement Learning in Computer-Use Agents}}{\textbf{arXiv, 2026}}
  \resumeProjectHeading
    {\href{https://papers.ssrn.com/sol3/papers.cfm?abstract_id=6727981}{Political Symbols and Urban Amenities: The Spatial Logic of Flag Placement in Istanbul}}{\textbf{SSRN, 2026}}
\resumeSubHeadingListEnd

%-----------RESEARCH EXPERIENCE-----------
\section{研究经历}
\resumeSubHeadingListStart

  \resumeSubheading
    {浙江大学 计算机科学与技术学院}{杭州 / 远程混合}
    {科研研究员；导师：朱霖潮 (Dr. Linchao Zhu)}{2025.02 -- 至今}

  \resumeSubSubheading
    {\textbf{CaptchaArena} $|$ 项目负责人 \& 一作 $|$ arXiv 2026}{2026.01 -- 至今}
  \resumeItemListStart
    \resumeItem{主导 \textbf{CaptchaArena}——首个大规模细粒度交互式 CAPTCHA 训练数据集：\textbf{5 万道谜题}、\textbf{20 类任务}、5 种交互模式，包含鼠标/键盘底层操作、浏览器可执行真值与 5 万条回放验证的截图—动作轨迹。}
    \resumeItem{构建数据集的推理监督管线，蒸馏 \textbf{4.6 万条多轮轨迹}：动作条件化的逐步视觉推理、像素级 grounding、独立 VLM 双重评审、强模型兜底与人工复核。}
    \resumeItem{训练 \textbf{CaptchaAgent}——统一 Qwen3.5-9B 计算机操作策略，覆盖全部 20 类任务且无任务专用适配器；基于 37,621 条逐轮 SFT 样本 + GRPO，Pass@1 从 11.4\% 提升至 70.5\% 再至 \textbf{71.7\%}（基座 $\rightarrow$ SFT $\rightarrow$ RL）。}
    \resumeItem{搭建分布式多模态 SFT/RL 基础设施：ZeRO-3/FSDP、异步浏览器 rollout、32K 上下文与可执行环境奖励，无需学习型奖励模型与人工奖励标注即可在线 GRPO；RL 进一步提升两个外部 CAPTCHA 基准上的迁移表现。}
  \resumeItemListEnd

  \resumeSubSubheading
    {\textbf{SeekJudge} $|$ 计算机操作智能体奖励建模 $|$ arXiv 2026}{2026.01 -- 2026.08}
  \resumeItemListStart
    \resumeItem{构建并人工标注 \textbf{300+ 条轨迹}的跨平台长程计算机操作智能体基准，覆盖网页、桌面与移动环境，提供轨迹级判定与稠密的步级监督信号。}
    \resumeItem{设计细粒度奖励标准，覆盖里程碑、正确决策、错误纠正、偏离路径行为与有害操作；开展人工评估与步级的模型—人类判定分歧分析。}
    \resumeItem{系统分析规则奖励与模型奖励信号的误差来源，识别多个 CUA 环境中规则与人类判定的分歧，论证长程智能体评估需要可学习的奖励模型。}
    \resumeItem{参与 \textbf{SeekJudge-9B} 的蒸馏与 SFT 训练，分析推理行为、输出格式与终止控制等失败模式，并迭代训练数据与监督策略。}
  \resumeItemListEnd

  \resumeSubSubheading
    {\textbf{智能体记忆 / 对话惯性 (Conversational Inertia)} $|$ ICML 2026}{2025.02 -- 2026.08}
  \resumeItemListStart
    \resumeItem{为长程 LLM 智能体实现任务定制的轨迹摘要，覆盖 Maze、WebShop、Wordle、SciWorld、TextCraft 等环境，在压缩多轮观察—动作历史的同时保留任务关键状态与进度。}
    \resumeItem{设计环境感知的记忆表示与摘要提示，选择性保留导航状态、已完成子目标、约束条件、历史动作与未解决信息，服务于后续智能体决策。}
    \resumeItem{在异构 \textbf{AgentGym} 环境上与多种上下文管理策略进行系统对比，分析历史信息保留与对话惯性之间的权衡；免训练摘要提示将 Maze 成功率从 85.5\% 提升至 \textbf{90.0\%}。}
    \resumeItem{对长上下文智能体开展失败与行为分析，构建注意力诊断链路（系统词元占比 $\rightarrow$ 系统注意力 $\rightarrow$ 任务成功率），并研究动作重复、探索行为与错误传播。}
  \resumeItemListEnd

\resumeSubHeadingListEnd

%-----------INDUSTRY EXPERIENCE-----------
\section{实习经历}
\resumeSubHeadingListStart

  \resumeSubheading
    {每日互动股份有限公司 (Merit Interactive Co., Ltd.)}{杭州}
    {AI 工程师实习生}{2026.01 -- 2026.03}
  \resumeItemListStart
    \resumeItem{设计端到端多模态检索系统，覆盖 \textbf{77 万+ 图片}，支持文搜图与图搜图，链路包含 FAISS ANN 检索、重排序与分数融合，向量召回延迟 \textbf{$<$50 ms}，端到端查询延迟约 100--500 ms。}
    \resumeItem{搭建覆盖 8 个主题域的多语言网页分类管线；人工标注约 6 千页面、整理 7.5 千+ 样本，系统对比 \textbf{13 种方案}（向量分类器、对比学习、LoRA、聚类、检索/重排序等）。}
    \resumeItem{将分类准确率从 62.5\% 提升至 \textbf{87.75\% Top-1}（97.6\% Top-2），并将最终 264 KB 轻量模型以容器化 REST API 形式部署上线。}
  \resumeItemListEnd

  \resumeSubheading
    {每日互动股份有限公司 (Merit Interactive Co., Ltd.)}{杭州}
    {算法实习生}{2024.07 -- 2024.08}
  \resumeItemListStart
    \resumeItem{基于原始用户日志用 SQL 构建移动应用分析数据集；分析使用、安装、卸载、用户分布与时间趋势，并自动化生成 Top-100 榜单报告。}
    \resumeItem{基于行为与时序特征，使用逻辑回归与 LightGBM 构建应用趋势预测模型，预测下周可能上榜的应用。}
  \resumeItemListEnd

\resumeSubHeadingListEnd

%-----------ADDITIONAL RESEARCH & TEACHING-----------
\section{其他研究与教学}
\resumeSubHeadingListStart

  \resumeSubheading
    {罗切斯特大学 Goergen 数据科学与人工智能研究院}{罗切斯特, NY}
    {DSCC 279 独立研究；导师：Prof. Cantay \c{C}al{\i}\c{s}kan}{2025.08 -- 2026.01}
  \resumeItemListStart
    \resumeItem{基于约 \textbf{74.8 万张 Google 街景全景图}构建地理空间视觉管线，包含图像清洗、坐标对齐、空间匹配、多尺度地理聚合与选区级特征构建。}
    \resumeItem{应用嵌套地理缓冲区、距离加权聚合与 Earth Mover's Distance，量化政治符号可见度在城市各区之间的结构性差异。}
  \resumeItemListEnd

  \resumeSubheading
    {罗切斯特大学 计算机科学系}{罗切斯特, NY}
    {本科助教，CSC 280: Computer Models and Limitations}{2025.08 -- 2025.12}

\resumeSubHeadingListEnd

%-----------SELECTED PROJECTS & AWARDS-----------
\section{竞赛与奖项}
\resumeSubHeadingListStart

  \resumeProjectHeading
    {\textbf{GIDS 生物医学数据科学黑客松} $|$ \textit{获奖}}{2026.08}
  \resumeItemListStart
    \resumeItem{\textit{MERFISH 空间转录组细胞类型分类：}主导建模管线，在小鼠脊髓数据（200 个基因、仅 5 千标注细胞）上预测 60 类细胞类型，基于表达、PCA 与空间/表达近邻标签特征构建种子 bagging 的 LightGBM 集成，引入生物学约束掩码与无泄漏交叉验证；在最终验证集排名中获\textbf{第 2 名}（共 30 队），公开榜单准确率 \textbf{80.7\%}（多数类基线 14\%）。}
  \resumeItemListEnd

  \resumeProjectHeading
    {\textbf{GIDS 生物医学数据科学黑客松} $|$ \textit{获奖}}{2025.08}
  \resumeItemListStart
    \resumeItem{\textit{Perturb-seq 基因表达预测：}开发增强版 Norman 方法管线，用于单/双基因扰动预测，结合稳健预处理、多模型回归与自适应系数估计；测试集 \textbf{RMSD 0.40}。}
  \resumeItemListEnd

  \resumeProjectHeading
    {\textbf{GIDS 生物医学数据科学黑客松} $|$ \textit{获奖}}{2024.08}
  \resumeItemListStart
    \resumeItem{\textit{单细胞多组学建模：}构建基于 MLP 的 peak-to-gene 预测管线，取得 \textbf{75\% 准确率}，优于线性回归与随机森林基线。}
  \resumeItemListEnd

  \resumeProjectHeading
    {\textbf{Finnov8 黑客松} $|$ \textit{Second Best Overall Award（总分第二名）}}{2024.03}
  \resumeItemListStart
    \resumeItem{在 4 人团队中取得 \textbf{0.67} 的个人最高分，通过特征工程、多输出 MLP 与 LSTM 模型将预测分数从 0.60 提升至 0.67。}
  \resumeItemListEnd

\resumeSubHeadingListEnd

%-----------LEADERSHIP-----------
\section{学生工作}
\resumeSubHeadingListStart
  \resumeProjectHeading
    {\textbf{宣传部部长}，罗切斯特大学中国学生学者联合会 (CSSA)}{2024.06 -- 2025.07}
  \resumeProjectHeading
    {\textbf{外联部成员}，罗切斯特大学中国学生学者联合会 (CSSA)}{2023.06 -- 2024.05}
\resumeSubHeadingListEnd

%-----------TECHNICAL SKILLS-----------
\section{技能}
\begin{itemize}[leftmargin=0.15in, label={}]
  \small{\item{
    {\sl 编程语言：} Python, C/C++, Java, SQL, JavaScript/HTML/CSS \\
    {\sl 机器学习/AI：} PyTorch, Hugging Face Transformers, Scikit-learn, FAISS, LoRA, SFT, GRPO, 多模态/VLM 训练 \\
    {\sl 系统/工具：} DeepSpeed ZeRO-3, Docker, Linux, Git, REST API, 分布式训练与检查点管理
  }}
\end{itemize}

\end{document}
